\documentclass[11pt]{article}
\usepackage{template}
\usepackage{amsmath,amssymb,microtype,tikz,graphicx,xurl}
\usetikzlibrary{arrows.meta,positioning,calc}
\setlength{\emergencystretch}{3em}
\raggedbottom
\SetKeys{date=2026-10-28,vol=015,title={每日有趣新闻｜时政、科学与数学},author={ChatGPT 编辑部},note={\small 采编窗口：2026年10月26—28日；截稿：北京时间10月28日上午。区分事件发生日期、报道日期和统计日期；预印本结论未经本刊独立证明核验。}}
\renewcommand{\newsitem}[3]{\par\medskip{\large\bfseries\color{newsblue}#1\par}{\footnotesize\color{newsgray}#2\par}\smallskip\color{black}#3\par\medskip{\color{newsgray}\rule{\linewidth}{0.25pt}}\par}
\newcommand{\newslead}[3]{\par\medskip{\LARGE\bfseries\color{newsblue}#1\par}\vspace{1mm}{\small\color{newsgray}#2\par}\smallskip\color{black}#3\par\medskip{\color{newsgray}\rule{\linewidth}{0.45pt}}\par}
\newcommand{\newssub}[1]{\par\smallskip\noindent{\bfseries\color{newsblue}#1}\quad}
\begin{document}
\begin{documents}
\raggedcolumns

\newslead{可再生能源可能首次不再刷新装机纪录：真正卡住增长的是电网}{能源深度｜IRENA总干事10月27日接受路透社采访；装机基准为2025年}{
过去五年，全球太阳能和风能新增装机屡创新高，以至于人们习惯把能源转型理解成一条只会向上的曲线。但国际可再生能源署（IRENA）总干事弗朗切斯科·拉卡梅拉10月27日在新加坡国际能源周表示，\textbf{2026年全球新增清洁能源装机可能低于2025年}，成为五年来首次放缓。他没有给出2026年具体预测值；把这番定性判断写成某个确定的新增吉瓦数，是超出证据的。\cite{ReutersGrid,IRENA2025}

\newssub{先看确定的基数。}IRENA今年4月发布的统计指出，2025年全球可再生能源装机总量达到5149吉瓦（GW），比上一年增加692吉瓦，增长15.5\%；新增电力装机中可再生能源约占85.6\%。装机容量是设备在规定条件下能够提供的额定功率，不是实际发出的电量。太阳能的夜间发电为零，风电输出随风况变化；因此5149 GW不能与火电站全年实际发电量直接相加比较。\cite{IRENA2025}
\begin{center}\begin{tikzpicture}[x=0.68cm,y=0.52cm,font=\scriptsize]
\draw[->,gray] (0,0)--(6.9,0);
\node[anchor=east] at (0,2.9) {2024};
\node[anchor=east] at (0,1.35) {2025};
\fill[newsblue!55] (0,2.55) rectangle (4.46,3.2);
\fill[newsblue!85] (0,1.0) rectangle (5.15,1.65);
\node[anchor=west] at (4.5,2.88) {4457 GW};
\node[anchor=west] at (5.2,1.32) {5149 GW};
\draw[thick,->,newsblue] (4.45,2.3)--(5.1,1.9);
\node[newsblue] at (3.5,1.85) {新增692 GW};
\node[align=center,gray,text width=6.6cm] at (3.3,-0.62) {全球可再生能源累计装机，2024值由2025总量减去当年净新增推算；来源：IRENA 2026年4月统计。2026年尚无本刊可核实的年度新增数。};
\end{tikzpicture}\end{center}\color{black}

\newssub{为什么发电设备越多，反而可能弃电？}电力系统至少包含发电、输电、配电、调度和用电。大型光伏电站可以在日照充足时同时达到输出高峰，但若远方城市负荷较低，或跨省输电线路、变压器、局部电网承载能力不足，调度系统就必须限制部分机组出力。这叫弃风弃光或发电量削减（curtailment）。它不是太阳能电池板坏了，而是某个时刻的电力无法以可接受的成本和安全条件送到需要它的地方。\cite{ReutersGrid}

拉卡梅拉认为，中国的可再生能源部署放缓是全球总量难以继续创纪录的重要原因，而澳大利亚、日本等国也面临电网扩容赶不上电源建设的问题。中国拥有庞大的新能源制造能力和装机规模，但\textbf{制造更多组件、安装更多机组与实际消纳更多电量，是三个不同的政策目标}。路透社引用他对中国弃电规模的跨国比较；本刊不把这一口径未经核对地换算成全国弃电率。\cite{ReutersGrid}

\newssub{储能能否代替电网？}电池储能可以把中午的太阳能搬到傍晚使用，抽水蓄能能够进行更大规模的时段转移；需求响应鼓励工业负荷在电力充裕时运行。它们都能降低弃电，却无法在所有地区、所有季节完全代替长距离输电。跨区域风光资源互补，需要更灵活的市场调度、联网和稳定的输电投资。电网建设本身也要处理环境、土地与居民利益的协调。

\newssub{放缓是否意味着能源转型失败？}不能如此推断。IRENA的判断针对\textbf{年度新增装机}，并不等于全球既有可再生能源总装机下降，更不等于当年发电量必然减少。2025年的692 GW还是异常高的比较基数。真正值得观察的是新增电网容量、输配电投资、弃电率、跨区交易、储能时长和最终发电量。能源转型的瓶颈正在从“能不能造出足够多的发电设备”，转向“能不能把电力系统作为整体改造”。
}

\newslead{局部停火开始执行：奥列什基平民救援能否跨过雷区与政治分歧？}{俄乌战争深度｜10月27日停火首日；路透社、美联社及《卫报》交叉核对}{
乌克兰赫尔松州第聂伯河东岸的奥列什基（Oleshky），自2022年起处于俄军控制之下，近年来遭战斗、洪水和交通中断多重打击。10月27日，一项在美国斡旋下达成的\textbf{局部人道停火}开始生效，目标是允许向城内运送食品、药品和饮水，并组织平民撤离。路透社称停火于当地时间凌晨3时开始，计划持续至11月1日。乌方人权监察官称首日尚未收到违反停火的报告；这是一项初步迹象，却不是对所有前线行动的独立实时监控。\cite{ReutersTruce,APTruce,GuardianTruce}

\newssub{为什么一座城需要单独谈判？}战线附近道路存在地雷和无人机威胁，车辆难以安全通过；长期断电、缺水、缺暖气使居民生活困难。多家媒体援引当地和乌方人士估计，城内仍有约1500至2000名平民，许多是高龄或行动不便者。不同当事方对人口规模估算并不一致，不能把这些数字写成已完成的人口普查。停火协议首先提供暂时不攻击特定救援活动的承诺，随后还需要排雷、车辆组织、路线确认、医护支持与居民登记。\cite{ReutersTruce,APTruce}

\newssub{最难的不是签字，而是安全通行。}报道显示，撤离车辆预计在停火开始数日后才能投入使用，因为道路需要排雷。另一项争议是居民撤往哪里：乌克兰总统泽连斯基担忧俄方可能将部分撤离者引向克里米亚或俄罗斯境内；俄方与乌方对通道控制的说法存在差异。居民是否能够自主选择目的地、是否允许独立人道机构接触，以及离开后是否有持续照护，比“多少辆大巴开进城”更能衡量行动质量。关于胁迫或筛查的具体指控仍需可独立核查的个案证据。\cite{APTruce,GuardianTruce}

\newssub{停火不等于和平。}局部停火只覆盖约定的区域和时段，并不要求交战双方停止其他战线的军事行动。10月27日，基辅桥梁及基础设施继续遭受俄军袭击；乌方公布了伤亡与损毁信息，部分细节仍待独立确认。与此同时，芬兰总统斯图布与挪威首相斯特勒同普京通话，试图推动降级并为拟议中的新一轮谈判争取欧洲参与。俄方仍强调其安全与领土要求，乌方则反对将俄占地区合法化。\cite{ReutersKyivBridge,ReutersNordic}

从外交逻辑看，能在一座城执行几天停火，至少证明敌对双方在有限议题上仍能协调；但它并不能证明在领土、安全保证、制裁或战后秩序等核心问题上已有共识。接下来几天最有价值的事实不是外交辞令，而是有多少居民在知情自愿下安全撤离、多少救援物资实际抵达、是否出现可证实的停火违约，以及停火结束后能否维持最低限度的人道通道。
}

\newsitem{2025年二氧化碳浓度再创新高：“排放”与“浓度”不是同一回事}{气候与公共政策｜世界气象组织10月27日发布第22期温室气体公报}{
世界气象组织（WMO）10月27日公布全球大气观测结果：2025年全球平均地表二氧化碳浓度为$426.8\pm0.2$ ppm，甲烷为$1948\pm2$ ppb，氧化亚氮为$339.1\pm0.1$ ppb，三者均创纪录。与1750年前的基线相比，它们分别相当于153\%、267\%和126\%。换言之，二氧化碳增加了53\%，甲烷增加了167\%，而不是“甲烷增加267\%”。\cite{WMOReport,WMOPress}
\begin{center}\begin{tikzpicture}[x=0.02cm,y=0.62cm,font=\scriptsize]
\draw[->,gray] (0,0)--(310,0);
\foreach \x in {100,200,300}{\draw[gray!30] (\x,0)--(\x,3.8);\node[gray] at (\x,-0.27) {\x};}
\fill[newsblue!75] (0,2.75) rectangle (153,3.27);
\fill[newsblue!55] (0,1.75) rectangle (267,2.27);
\fill[newsblue!35] (0,.75) rectangle (126,1.27);
\node[anchor=east] at (-4,3) {CO$_2$};\node[anchor=east] at (-4,2) {CH$_4$};\node[anchor=east] at (-4,1) {N$_2$O};
\node[anchor=west] at (156,3) {153\%};\node[anchor=west] at (270,2) {267\%};\node[anchor=west] at (129,1) {126\%};
\node[align=center,gray,text width=6.6cm] at (150,-0.95) {相对于1750年前浓度基线（基线=100\%），不是年度排放增长率。来源：WMO《温室气体公报》第22期，2025年观测。};
\end{tikzpicture}\end{center}\color{black}
\textbf{排放量}是一定时间内向大气释放的气体质量，\textbf{浓度}是大气里已经积累了多少。即使年度排放开始下降，只要排放仍超过海洋、陆地生态系统等吸收的总量，浓度仍可能上升。因此某一年浓度创新高，不足以单独判断当年排放增速。WMO报告称2024至2025年的二氧化碳浓度增量小于上一年度异常偏高的增量，但仍略高于过去十年的平均水平。\cite{WMOReport}
WMO同时提醒，厄尔尼诺相关干旱可能削弱陆地碳汇，并提到电力设备中使用的六氟化硫（SF$_6$）及氢气泄漏带来的间接气候效应。六氟化硫排放量相对较小，却具有很高的单位质量增温潜势；氢气本身不是传统的长寿命温室气体，但可能通过大气化学反应影响甲烷寿命。把这些风险纳入能源基础设施设计，并不是反对清洁能源，而是要求在转型过程中同时管理新的排放源。\cite{WMOPress,ReutersWMO}
}

\newsitem{Oracle云引入Google TPU：AI算力市场正在分化}{人工智能产业｜10月27日Oracle、Google官方公告及路透社报道}{
Oracle与Google10月27日宣布，Google张量处理器（TPU）将通过Oracle Cloud Infrastructure（OCI）提供服务。Oracle称这是Google以外首个提供其TPU的超大规模云平台；同时Oracle扩大与NVIDIA Vera Rubin系统合作，目标为“多吉瓦级”部署，并计划提供Cerebras、SambaNova等公司的推理芯片。官方宣布的是产品可用性与扩张计划，并不意味着所有规划算力已完成建设、客户已购买或未来收入已实现。\cite{OracleTPU,OracleNV,ReutersOracle}
训练大型模型通常需要高带宽互连、巨量矩阵运算与复杂软件栈；部署模型时的推理任务则在响应延迟、吞吐量、显存容量、能耗和成本之间作不同取舍。不同客户可能偏好NVIDIA GPU、Google TPU或专门面向推理的加速器。Oracle在部分NVIDIA系统中还计划采用自研Acceleron网络芯片，说明竞争不仅发生在处理器本身，也发生在网络、内存、调度和软件生态之间。\cite{ReutersOracle}
路透社指出，Oracle维持2027财年约900亿至950亿美元资本开支计划，但没有披露这次新增芯片采购的完整融资结构。NVIDIA管理层曾以每吉瓦最高约400亿美元描述特定系统的潜在收入，这是一种产业估算，不能乘以“多吉瓦”就当成Oracle已签署的合同金额。宣布算力规模、交付算力、客户使用率和最终利润仍是四个不同层次。\cite{ReutersOracle}
}

\newsitem{地球辐射带的“不可穿透墙”并非绝对：日本卫星发现本地加速}{空间物理｜京都大学10月27日发布；研究于10月26日发表于GRL}{
地球磁场能够俘获高能带电粒子，形成范艾伦辐射带。以往观测显示，在地球周围某些较低高度区域，超高能电子很难长期进入，因此研究者把相应边界称为“不可穿透墙”。京都大学、名古屋大学等团队分析日本“荒濑”（Arase/ERG）卫星在2024年5月强磁暴期间的数据后，报告在边界内侧发现多兆电子伏电子的局部增强，研究于10月26日刊登在《地球物理研究快报》。\cite{KyotoArase,MiyoshiGRL}
\begin{center}\begin{tikzpicture}[scale=.73,font=\scriptsize]
\shade[inner color=newsblue!6,outer color=newsblue!20] (0,0) circle (2.05);
\draw[newsblue!70,thick] (0,0) circle (2.05);
\draw[red!65!black,dashed,very thick] (0,0) circle (1.35);
\fill[blue!45!green!35] (0,0) circle (.73);
\fill[newsblue!80] (0,0) circle (.38);
\node[white] at (0,0) {地球};
\draw[->,very thick,green!55!black] (1.72,1.1) .. controls (1.35,.5) and (1.45,-.4) .. (.96,-.85);
\node[green!40!black,align=left,anchor=west] at (2.2,1.25) {合唱波与局部加速\\（概念示意）};
\node[red!65!black,anchor=west] at (2.1,-.45) {经验性屏障};
\draw[->,red!65!black] (2.0,-.45)--(1.26,-.45);
\node[gray,align=center,text width=6.8cm] at (0,-2.8) {非比例轨道图；电子可在边界内侧被加速，并非穿透实体墙。依据：京都大学2026年10月27日研究说明。};
\end{tikzpicture}\end{center}\color{black}
关键证据不只是粒子数量增加：卫星还在相近地点和时刻观测到强烈的\textbf{合唱波}（chorus waves），以及异常稀薄的背景等离子体。团队提出，磁暴使等离子体层收缩，电磁波在特殊环境下持续把部分电子加速至更高能量。这些电子不一定从外侧直接穿过原本的边界，而可能\textbf{在边界内侧被加速产生}。研究将观测与物理解释联系起来，但单次极端事件不能不加限定地外推到所有磁暴。\cite{KyotoArase}
结果对卫星电子设备的辐射风险评估有现实意义：过去被视为相对稳定的轨道区域，在极端空间天气下也可能出现新的高能粒子环境。所谓“墙”是经验性的动力学屏障，不是固体结构，更不是地球磁场突然破了实体洞。
}

\newsitem{数学拾趣｜一棵树上，最少点亮几个顶点才能点亮整棵树？}{组合数学｜arXiv:2610.35745；10月26日提交，27日进入新列表；预印本}{
设一棵树的顶点起初只有部分被标记为“活跃”。给定阈值比例$\rho\in[0,1]$，一个尚未活跃的顶点$u$，只要至少有$\lceil\rho d(u)\rceil$个邻居已活跃，就在下一轮变为活跃；这里$d(u)$是顶点度数，尖括号表示向上取整。已经活跃的顶点不会恢复。若从某个初始集合出发，经过有限轮能使全图活跃，就称它为\textbf{传染集合}。这是一种图论的阈值扩散模型，不是在预测真实传染病。\cite{ArxivContagion}
最简单的例子是一条路径：若$\rho\leq\tfrac12$，从端点点亮一个顶点，活跃状态就能沿路径逐个传播；若$\rho=1$，每个内部顶点要求全部邻居已活跃，扩散规则更苛刻。Jeong-Ok Choi与Collin McDevitt的新预印本研究固定顶点数$n$时，\textbf{所有树中“最小传染集合的大小”最大可以有多大}。注意先对每棵树取最小值，再在所有$n$顶点树中取最大值，量词顺序不能交换。\cite{ArxivContagion}
作者声称得到完全$d$叉树上的精确值，并确定上述极值；方法包括寻找互不相交的抗扩散集合，每个集合都必须与任意传染集合相交，从而给出下界，再构造达到边界的树。作者还证明在$0<\rho\leq1$、$n\ge2$范围内，从1到最大值之间每个整数都能由某棵树实现。论文尚未经本刊独立逐行核验，也没有本刊确认的Lean内核证书。\cite{ArxivContagion}
}

\newsitem{数学短讯｜零和序列两个经典结构猜想遭到更大范围反例挑战}{数论与组合｜arXiv:2610.35819；10月26日提交}{
Christian Elsholtz在新预印本中研究有限阿贝尔群$C_n^r$上的零和序列问题。这里$C_n$是$n$阶循环群，$C_n^r$是其$r$个直积；序列中的元素允许重复。Erd\H{o}s--Ginzburg--Ziv常数$s(C_n^r)$衡量强制出现恰有$n$项和为零的子序列所需的最小长度。性质C、D猜想试图刻画极端地避免零和子序列的序列必须具有什么重复结构。\cite{ArxivZeroSum}
论文声称只要$n\equiv2\pmod4$、$n\ge6$且$r\ge4$，两项结构性质均不成立，并给出$s(C_n^r)\ge2^r(n-1)+2^{r-2}+1$的下界。作者把此前$C_6^r$的反例推广到整个同余类。这是对结构猜想的反例主张，并非求出了所有参数下零和常数的精确值；尚未经本刊独立验证或同行评议。\cite{ArxivZeroSum}
}

\newsitem{中国AI概念股回调：股价缩水与技术能力下降不是同一命题}{中国市场与科技产业｜路透社10月26日专题；10月27日市场续跌}{
中国内地科技股在2026年夏季经历大幅上涨后转入调整。路透社10月26日市场分析指出，科创50指数在不到四个月内回撤约40\%，涉及市值损失估计约4800亿美元；沪深300指数相较6月高点也下跌约20\%。不同指数的成分股、估值基准和计算窗口不同，不能把某一指数跌幅解释成整个中国科技产业价值蒸发了同样比例。\cite{ReutersChinaAI,GuardianChinaAI}
AI芯片、光通信、机器人和模型企业此前承载了较高增长预期。一些公司营收仍在增长，但投资者重新评估盈利兑现速度、市场竞争、出口限制和融资成本。即便基本面继续改善，只要未来利润预期从“极其乐观”降为“比较乐观”，估值也可能大幅回撤。政策支持关键技术研发，不等于替每家上市企业保证利润。\textbf{技术自主可控的国家目标与具体股票的投资回报，是两个不同层次的评价。}
本期不把“泡沫破裂”作为客观证实的产业结论：它是部分媒体和市场人士对估值变化的描述。判断产业进展仍需查看产品性能、交付规模、收入质量、研发成果和长期成本结构。\cite{ReutersChinaAI}
}

\newsitem{美国拟放宽油气甲烷泄漏规则：降低合规成本与增加排放的政策权衡}{美国公共政策｜10月27日政府表态；美联社与《华盛顿邮报》报道}{
美国政府10月27日宣布计划放宽针对油气开采设施的甲烷泄漏监管，涉及小型油井、常规放空燃烧限制，以及由第三方遥感发现重大泄漏的“超级排放源计划”。政府及环保署负责人主张，现有规则增加小型生产商合规负担，放松监管有助于降低成本、扩大国内能源供应。环保团体则认为，减少泄漏检测和维修义务会让本可回收利用的天然气排入大气，并带来气候与公共卫生损害。\cite{APMethane,WPMethane}
必须区分\textbf{提出放宽规则}与\textbf{新规则已经完成法定程序并生效}。报道的是政策方向与拟议修订，具体法律文本、公众意见征集和司法挑战可能影响实施。美国环保署现有超级排放源计划网页显示，此前规定的实施时间已延至2027年1月22日；该网页可核查原制度，但不能单独证明10月27日拟议修订的全部条款。\cite{EPASuper}
甲烷是天然气的主要成分，短期增温效应强于二氧化碳；泄漏同时意味着能源损失。监管难点在于，某些边际小井检测成本相对于产量较高，但部分低产井又可能贡献不成比例的泄漏。合理评估应公开监测数据、实际减排量、合规费用及受影响社区的健康指标，而不是只用“减负”或“环保”口号代替分析。\cite{WPMethane}
}

\newsitem{欧盟要求按更高升温情景规划城市：适应气候不是放弃减排}{欧洲公共政策｜欧盟委员会10月27日发布气候韧性框架}{
欧盟委员会10月27日提出气候韧性框架，要求成员国把热浪、干旱、洪水、森林火灾和海平面变化纳入基础设施、城市与公共财政规划。委员会估计，到2050年，欧洲适应未来气候条件每年需要约700亿欧元投资。欧洲投资银行计划在2026至2030年期间为气候适应配置至少300亿欧元融资。\cite{EUClimate,ANSAClimate}
这里的700亿欧元是\textbf{投资需求估算}，并非欧盟已通过的年度预算拨款。新的框架仍需在欧盟制度内推进，成员国实际规划与财政承诺也待观察。应对措施可能包括建筑遮阳与隔热、洪泛区管理、城市绿地、预警系统和公共卫生响应。\cite{EUClimate}
适应与减排解决不同时间尺度的问题：即使今天降低排放，已经积累的大气温室气体仍会在未来多年影响气候，基础设施必须为不可避免的风险做好准备；但若只投入防洪、空调和灾后修复而不继续减排，未来需要适应的风险会不断增加。采用“3摄氏度升温情景”做压力测试，并不是政府宣布一定要让地球升温3摄氏度。\cite{FTClimate}
}

\par\medskip{\large\bfseries\color{newsblue}参考资料与编辑说明}\par
{\scriptsize 主要材料截至北京时间2026年10月28日上午。政治报道区分当事方声明与独立核验；WMO数字为2025年观测结果；数学论文为2026年10月26日提交的预印本。三张插图均为本刊原创TikZ信息图或概念示意。}\par
\begingroup\fontsize{6}{7.1}\selectfont
\renewcommand{\refname}{参考文献}
\bibliographystyle{plainurl}\bibliography{refs}
\endgroup
\end{documents}
\end{document}